%%%%%%%%%%%%%%%%%%%%%%%%%%%%%%%%%%%%%%%%%%%%%%%%%%%
%% LaTeX book template                           %%
%% Author:  Amber Jain (http://amberj.devio.us/) %%
%% License: ISC license                          %%
%%%%%%%%%%%%%%%%%%%%%%%%%%%%%%%%%%%%%%%%%%%%%%%%%%%

\documentclass[a4paper,11pt]{book}
\usepackage[T1]{fontenc}
\usepackage[utf8]{inputenc}
\usepackage{lmodern}
%%%%%%%%%%%%%%%%%%%%%%%%%%%%%%%%%%%%%%%%%%%%%%%%%%%%%%%%%
% Source: http://en.wikibooks.org/wiki/LaTeX/Hyperlinks %
%%%%%%%%%%%%%%%%%%%%%%%%%%%%%%%%%%%%%%%%%%%%%%%%%%%%%%%%%
\usepackage{hyperref}
\usepackage{graphicx}
\usepackage[english]{babel}

%%%%%%%%%%%%%%%%%%%%%%%%%%%%%%%%%%%%%%%%%%%%%%%%%%%%%%%%%%%%%%%%%%%%%%%%%%%%%%%%
% 'dedication' environment: To add a dedication paragraph at the start of book %
% Source: http://www.tug.org/pipermail/texhax/2010-June/015184.html            %
%%%%%%%%%%%%%%%%%%%%%%%%%%%%%%%%%%%%%%%%%%%%%%%%%%%%%%%%%%%%%%%%%%%%%%%%%%%%%%%%
\newenvironment{dedication}
{
   \cleardoublepage
   \thispagestyle{empty}
   \vspace*{\stretch{1}}
   \hfill\begin{minipage}[t]{0.66\textwidth}
   \raggedright
}
{
   \end{minipage}
   \vspace*{\stretch{3}}
   \clearpage
}

%%%%%%%%%%%%%%%%%%%%%%%%%%%%%%%%%%%%%%%%%%%%%%%%
% Chapter quote at the start of chapter        %
% Source: http://tex.stackexchange.com/a/53380 %
%%%%%%%%%%%%%%%%%%%%%%%%%%%%%%%%%%%%%%%%%%%%%%%%
\makeatletter
\renewcommand{\@chapapp}{}% Not necessary...
\newenvironment{chapquote}[2][2em]
  {\setlength{\@tempdima}{#1}%
   \def\chapquote@author{#2}%
   \parshape 1 \@tempdima \dimexpr\textwidth-2\@tempdima\relax%
   \itshape}
  {\par\normalfont\hfill--\ \chapquote@author\hspace*{\@tempdima}\par\bigskip}
\makeatother

%%%%%%%%%%%%%%%%%%%%%%%%%%%%%%%%%%%%%%%%%%%%%%%%%%%
% First page of book which contains 'stuff' like: %
%  - Book title, subtitle                         %
%  - Book author name                             %
%%%%%%%%%%%%%%%%%%%%%%%%%%%%%%%%%%%%%%%%%%%%%%%%%%%

% Title and Subtitle
\title{\Huge \textbf Absurd Grace\\ \small \textit{A Field Guide For The Culturally Overloaded and Chronically Online}}


% Author
\author{\textsc{Michael Hacker}}


\begin{document}

\frontmatter
\maketitle

%%%%%%%%%%%%%%%%%%%%%%%%%%%%%%%%%%%%%%%%%%%%%%%%%%%%%%%%%%%%%%%
% Add a dedication paragraph to dedicate your book to someone %
%%%%%%%%%%%%%%%%%%%%%%%%%%%%%%%%%%%%%%%%%%%%%%%%%%%%%%%%%%%%%%%
\begin{dedication}
For those who find the humor in it all.
\end{dedication}


\mainmatter
\tableofcontents

%%%%%%%%%%%
% Preface %
%%%%%%%%%%%
\chapter*{Preface}
\textbf{\\Disclaimer\\}

This is a field guide, not a manifesto. There are no blueprints here, no secret keys, no “10 Proven Hacks for Kingdom Influence.” What you’ll find instead is a scattered constellation of observations, provocations, laments, and softly spoken prayers—each written with the humble awareness that grace is not a system, but a scandal. A mystery. A miracle.

Beyond here are ideas that may challenge your comfort, unsettle your assumptions, or even (God forbid) make you laugh mid-repentance. Your experience is your own—handle it with care and with coffee. The author takes no responsibility for spiritual awakenings, existential spirals, or sudden bursts of gratitude in traffic.

The views presented here do not constitute legal advice, denominational decrees, or liturgical dress codes. This is Kierkegaard with a Wi-Fi connection. This is Leibniz after a long fast. The thoughts within are ancient, absurd, and—unfortunately for the algorithm—not optimized for relevance. They have survived emperors, empires, and engagement metrics.

This guide is offered not as gospel, but as graffiti on the walls of the modern temple—an invitation to think differently, to live more deeply, and to question the sacred branding strategies masquerading as truth. It is written for a world galloping toward technological transcendence, where digital gods demand our constant devotion. Against that tide, this book prays (and protests) for a different kingdom.

At its core, this work is shaped by a prayer—a longing not for performance but for presence. A lament that doubles as praise. A cry from the depths of a digitally saturated soul that still dares to believe in the riches of divine provision. It is a field guide for those who believe that grace is not earned, not intellectualized, and certainly not withheld by gatekeepers in trendy sneakers.

We live, not to prove our worth, but because we have been overwhelmed by a love that renders all performance obsolete. That is the real disclaimer: once grace takes root, there is no going back.

\textbf{\\Introduction\\}
Let’s be honest—grace doesn’t make sense. It’s not efficient. It doesn’t monetize. It doesn’t scale. There’s no subscription model, no branded content strategy, no influencer collab that can fully explain what grace is or why it keeps showing up like a divine party crasher. Grace is the worst business plan in human history—and yet, it’s the only thing that actually changes anything.

This book is about that kind of grace. The kind that walks into your disaster of a life without knocking and says, “Hey, I brought snacks.” The kind that doesn't ask for your résumé or your sanctification status. It just sits down, hands you a cold drink, and tells you you're still wanted.

But grace isn’t just a warm hug or a theological buzzword. It’s a force, not a badge. It doesn’t exist to be admired like a museum piece or dissected like a frog in Sunday school. Grace moves. And when it hits you, it doesn’t just pat you on the back—it recruits you. It says: Trust. Give. Follow. Love. Not because you have to, but because you finally can.

To suggest that grace somehow makes giving or trusting less important is to commit a kind of theological category error. It’s like confusing the music for the dance, or worse, acting like the dance floor is optional. But grace without trust is sentiment. Grace without generosity is just a lecture. Grace that doesn’t lead to participation is grace that hasn’t actually landed.

And that’s the real scandal: not just that grace forgives, but that it transforms. It calls you off the couch, out of the feed, and into the messy, embodied life of discipleship. It pulls your hands off your wallet, your ego off its throne, and your eyes off the algorithm. It teaches you to live again.

This book is a field guide for that journey—written in the language of memes, ministry, burnout, and late-night doubt. You’ll meet prophets in sneakers, gatekeepers with broadband, and revivalists who’ve trademarked the Holy Spirit™. You’ll laugh, roll your eyes, and maybe even throw the book once or twice. But somewhere between the satire and the sacred, I hope you’ll feel it:

That absurd, unscalable grace.

Not asking for anything from you—but pulling you into something bigger anyway.

Because the truest things in life can’t be bought, branded, or streamed.
They can only be received.

And once received—lived.

\chapter*{Prologue}

Prologue
Let’s begin with a confession: this book should not exist.

Not because it’s heretical (though some may think so), nor because it strays too far from orthodoxy (though others may claim that too), but because what it attempts to capture defies articulation. Grace, in its truest form, is too vast to be systematized, too sacred to be commodified, and far too absurd to be marketed in a world obsessed with productivity, prestige, and platforms.

And yet here we are.

This is a book about grace in the modern world—a world where preachers have YouTube intros, prophets sell merchandise, and the sacred is often drowned in SEO optimization. A world where we’ve confused going viral with being anointed and where the Holy Spirit is expected to show up between Wi-Fi outages and worship sets.

The thesis of this work is simple: grace is absurd. It doesn’t play by the rules. It shows up in the wrong places. It blesses the wrong people. It forgives what should be unforgettable. And it refuses to stay inside the clean, curated boundaries of institutional religion or digital ministry.

Grace Is Not a Brand Strategy
In an era where every soul is a metric and every sermon a reel, we’ve watched the gospel become oddly symmetrical with the influencer economy. Some pulpits now operate like press conferences. Preachers enter to intro music. And the message? It often comes dressed in TED Talk polish, armed with buzzwords, and aimed not at the brokenhearted, but at the algorithm.

Here, grace is no longer a divine ambush. It's a ladder—climbed through insight, aesthetics, and social utility. The implicit message? Prove you “get it.” Perform it persuasively. Only then do you belong.

But the tragedy runs deeper. Many of those preaching this message—this polished, persuasive gospel of spiritual merit—have not themselves been wrecked by grace. Not really. Because real grace, the disruptive kind, doesn’t leave room for pretense. It unseats the ego, undoes the self-made, and ruins any theology that makes a brand out of blessing.

It is often the poor, the outcast, the overlooked—the ones without platforms—who know grace best. They’ve lived it. Not studied it. Not packaged it. Received it. And when these are the very people being told they must now earn their place in the assembly by proving their grasp on what they’ve already lived? That’s not just wrong—it’s alienating.

This isn’t Levitical stewardship. It’s pharisaical gatekeeping. It is intellectualism parading as holiness. It is the absurd turned upside down—requiring clean resumes before administering rescue. But grace doesn’t require a résumé. It writes a new story.

Grace doesn’t climb on anyone’s shoulders. It kneels. It washes feet. It disappears into the crowd, not to hide, but to serve without needing a spotlight. And when it shows up, it doesn’t wait for approval. It just gets to work.

Four Movements, One Story
Absurd Grace is divided into four parts.
Faith in the Digital Spotlight
In a world dominated by curated feeds and instant visibility, faith often finds itself shaped by the architecture of platforms rather than the architecture of the soul. The digital spotlight demands that everything be seen, shared, and monetized—and spiritual life is no exception.

Under this glare, sermons become soundbites, prayer becomes content, and sacred rhythms are flattened into marketable moments. While digital tools can amplify the message of the gospel, they can also distort it—encouraging performance over presence and influence over intimacy with God.

Still, true faith resists reduction. It thrives in silence, in anonymity, in obedience away from the camera lens. It invites believers to carry a light that shines not through popularity, but through quiet faithfulness. In the digital age, the task is to be both present in the culture and yet untouched by its hunger for spectacle.

Faith in the Age of Influencers
The modern spiritual landscape is increasingly shaped by influencers—figures who command platforms, audiences, and curated personas. In this climate, visibility often substitutes for virtue, and style can easily eclipse substance.

Influencer culture rewards charisma, but biblical faithfulness requires character. Followers can be counted, but can they be discipled? When sermons are measured by engagement metrics and theology is simplified for the algorithm, spiritual depth suffers.

True ministry, however, isn’t a brand. It’s a burden. A calling. A cross. The life of Christ points toward obscurity over celebrity, self-giving over self-promotion. In an age of likes and shares, faith calls for a deeper witness—a way of life marked by humility, love, and the slow work of transformation.

Reclaiming Real Life in a Virtual World
We are more connected than ever and yet often feel more distant—from one another, from our own bodies, and from God. Screens mediate our relationships and interrupt our presence. Even worship is increasingly consumed rather than participated in.

Christian faith is deeply incarnational. It affirms the importance of flesh, of space, of time. The sacraments are not digital. The fellowship of believers is not virtual. The Body of Christ is not just a metaphor—it is embodied, present, gathered.

To reclaim real life is to return to simplicity and slowness. It is to light a candle rather than swipe. To show up rather than log in. To build community that isn’t curated. In doing so, we honor the God who did not stay distant, but took on flesh and dwelt among us.

The Sacred Struggle for Authentic Ministry
Ministry in the digital era is a sacred struggle. The pull toward self-promotion, marketability, and platform is constant. The temptation is to exchange depth for reach, truth for trend, and calling for content.

But authentic ministry is not built on algorithms—it is built on altars. It is formed in the quiet hours of prayer, the tear-streaked faces of pastoral care, and the fidelity of teaching truth when it costs something. It is not flashy. It is faithful.

This struggle is holy. It is the refusal to allow the sacred to be sold. It is the stewardship of mystery in a world that commodifies meaning. It is, as Kierkegaard might say, the willingness to be a "witness to the truth" in a culture that craves entertainment.

In this, absurd grace remains the answer. Not as a product, but as a presence. Not as a performance, but as a person—Jesus Christ, who is the same yesterday, today, and forever.

"Be not conformed to this world: but be ye transformed by the renewing of your mind..." — Romans 12:2, KJV

Each chapter is a mini-book, a standalone meditation, sermon, and satire. Together, they form a mosaic of sacred subversion.

A Bit About Tone
This book will not always be solemn. Grace isn’t solemn. Sometimes it’s funny. Often it’s inconvenient. Grace slips in through the back door while we’re busy rehearsing our public testimonies. It’s not always polite. Sometimes it crashes the party and breaks the fine china.

So yes, you will find humor here. Satire, even. Because sometimes the only sane response to a deeply broken system is to laugh at it—then dismantle it with holy irreverence.

But you will also find a sincere ache. A yearning for what is real, what is holy, what is not for sale. This book is equal parts lament and love letter.

A Note on Scripture
Biblical references throughout this book are drawn from the King James Version (KJV)—not out of nostalgia or elitism, but because the poetry of the KJV captures the gravity and grace of Scripture in a way that modern translations sometimes flatten. When Scripture speaks in this book, let it do so with the thunder and tenderness of the old tongue:

"But God commendeth his love toward us, in that, while we were yet sinners, Christ died for us." — Romans 5:8, KJV

What could be more absurd than that?

A Quiet Rebellion
The spirit of this book is not rage, but reverence. Not outrage, but invitation. We do not rage against the modern world—we simply choose to live differently within it. To turn down the volume. To reject clickbait theology. To stop performing our faith and start practicing it.

If this book has a central image, it is this: a single candle burning in a dark, noisy room. That’s grace. Not a spotlight. Not a laser show. A flickering light that refuses to be snuffed out.

And if this book has a prayer, it is this:

That you would come to know this absurd grace. That you would learn to receive it, even when it makes no sense. That you would find in it the courage to love God, neighbor, and self—with all your heart, mind, soul, and strength. And that you would carry that love quietly, faithfully, without applause—

Even if no one ever follows you for it.

Welcome to Absurd Grace.

\chapter{Faith in the Digital Spotlight}

In the age of digital media, where every post has the potential to go viral and influence is often measured by clicks, likes, and shares, the intersection of religion and media has become more complex than ever. The themes explored in this first section, from the commercialization of faith to the role of influencers in spirituality, challenge the traditional notions of ministry, responsibility, and sacred duty.

\section{The Gospel According to Big Mike}

In a culture obsessed with fame, image, and algorithmic validation, Chapter 1 critiques how modern society repackages greatness into consumable icons—figures like Michael Jordan, Michael Jackson, and celebrity preachers—transforming the divine imago Dei into a marketable brand. Drawing on voices like Kierkegaard, Spurgeon, and Kimmerer, the chapter explores how this commodification distorts our understanding of faith, identity, and the nature of God.

The chapter satirizes influencer culture with love and theological depth, arguing that grace stands in radical opposition to the logic of branding. Grace, unlike curated personas, cannot be optimized or sold. Instead, it invites a quieter, deeper kind of greatness—one rooted in presence, humility, and a return to embodied faith. Through this lens, the chapter calls readers to exit the matrix of performance and embrace the absurd beauty of unmarketable grace.

\section{Holy LOL: The Meme-ification of Jesus}

The concept of Jesus as a meme challenges traditional notions of reverence and holiness. The rapid spread of humor and irreverence in modern social media has made it difficult to maintain sacred spaces, and in the process, divinity has been reduced to a punchline. The sacrilege of turning the Savior of the world into a meme seems antithetical to Christian belief, but it also reveals a deeper cultural phenomenon: the postmodern desire to defang sacred concepts.

Kurt Gödel’s incompleteness theorems offer an intriguing lens through which to view this issue. Just as Gödel proved that not all truths could be logically derived, we might argue that not all spiritual truths can or should be distilled into easily consumable pieces of entertainment. However, this trend reveals a cultural need for levity in a time of existential and social uncertainty. 

The question then becomes, how do we find the balance between reverence and relatability? While Kierkegaard argued that a leap of faith is required to grasp spiritual truths, modern Christians may find themselves compelled to laugh in the face of the unexplainable, as the pressures of the digital world distort sacred narratives into meme fodder.

\section{Prophets, Podcasters, and Other End-Time Occupations}

The proliferation of digital media has created a new breed of prophets—one whose words spread not through solemn sermons but through podcasts, live-streams, and social media posts. These digital prophets offer insights and predictions about everything from the apocalypse to personal success. But in a world oversaturated with content, who is truly called to speak on behalf of God, and who is merely a content creator trying to build a following?

Drawing on Immanuel Kant’s ideas of moral imperative, we can begin to question whether these new digital prophets are living up to their ethical responsibilities. Kant argued that we must treat individuals as ends in themselves, not merely as means to an end. This is a critical challenge for influencers and digital pastors who may see their followers as commodities or stepping stones to fame rather than souls to care for.

Jesus’ teachings on humility and selflessness stand in stark contrast to the self-promotion so common in the modern influencer economy. For those who claim to be prophetic, the true test lies not in their follower count but in their commitment to serving others with integrity and sincerity.

\section{The Gatekeeper’s Lament: Ministry in the Age of Monetization}

Ministry has always been a vocation of sacred responsibility, yet the rise of digital platforms has made it increasingly difficult to maintain the integrity of the call. The monetization of ministry, where leaders are incentivized by clicks, shares, and online donations, presents a challenge for those called to the work of the gospel. 

Drawing inspiration from the Levitical code, which emphasized the sanctity of sacred work, it becomes clear that ministry is not a profession to be treated as a business transaction. Leibniz’s concept of pre-established harmony can be applied here: just as the cosmos operates in a divine order, so too should ministry maintain its original purpose, focused on service rather than self-interest.

The digital age has turned ministers into influencers, with algorithms determining which voices rise to the top. The allure of fame and wealth that comes with widespread recognition can be tempting, but it threatens to undermine the core of Christian ministry, which is to serve, love, and teach with humility, rather than sell a product.

\textbf{\\Part I Summary\\}

In this section, we explored the troubling intersection of media, celebrity, and faith. From the commodification of the gospel to the rising tide of digital prophets, these chapters challenge us to rethink what it means to faithfully serve in an age where visibility often trumps truth. In the next section, we will continue to explore the impacts of digital culture on faith, moving from the complexities of influence to the profound implications of our online lives.

\chapter{The Gospel According to Big Mike}

It begins, as all modern myths do, with a montage.

The camera pans across the faces of greatness: Michael Jordan levitating above defenders, tongue out, Nike-clad feet defying physics. Cut to Michael Jackson moonwalking across the stage, sequined glove catching the spotlight. Then a montage of mega-preachers with million-dollar smiles, bathed in LED lights and fog machines, speaking in hashtags and sermon series titles that sound suspiciously like TED Talks.

Enter the age of Big Mike.

The Algorithm of Glory

Greatness, in this world, is no longer a calling—it’s a campaign. And Big Mike isn’t a person; he’s a prototype. Jordan. Jackson. That other Big Mike with the mic and the merch. Together they form the Holy Trinity of Influence: talent, image, and ubiquity. What once required prophets now only demands production value.

We live under the rule of the curated self. But this isn’t just a cultural footnote; it’s a theological crisis. In a world where identity is algorithmically reinforced, what room is there for the unquantifiable grace of God? We’ve gone from imago Dei to imago Brand.

As Søren Kierkegaard might say—were he tweeting from a blue-check account—we are a generation in despair, not because we have too little, but because we are too much: too seen, too staged, too sold. He called it the sickness unto death. We might call it the sickness unto likes.

Manufactured Transcendence

Grace, by contrast, refuses to go viral. It doesn’t fit the brand strategy. It cannot be quantified, monetized, or streamlined into a YouTube ad campaign. Charles Spurgeon didn’t need a drone shot or a soundboard—he had tears, a cracked voice, and the Word. The Spirit doesn’t descend through Wi-Fi. It groans through the cracked ceilings of humility.

The world hands us idols—polished and palatable. We bow, not before golden calves, but golden brands. And while Leibniz tried to justify the ways of God with elegant logic, today we justify the ways of God by engagement metrics and influencer partnerships. But the calculus of heaven cannot be computed in impressions. It’s more absurd than that. It’s grace.

A Box Made of Mirrors

We’re not just trapped in the matrix; we’re subscribing to it monthly. The great reversal has occurred. Up is down. The first are last. And greatness now comes shrink-wrapped and endorsed by sneaker lines. The danger isn’t that we admire excellence—it’s that we mistake spectacle for substance. As Kant would warn, we’ve replaced duty with dopamine.

But the gospel doesn’t play by these rules. It shows up in manger hay, not prime time. It rides donkeys, not limousines. It is foolishness to the Greeks, scandal to the Jews, and shadowbanned by modernity.

Exit the Brand

Grace is the last free thing. And it begins when we unplug—not just from devices, but from delusions. It begins when we stop scrolling and start seeing. When we stop performing and start confessing. When we stop asking, “How can I be followed?” and start asking, “Who am I following?”

Robin Wall Kimmerer reminds us that real knowing is relational. You cannot commodify a tree. You cannot trend a sacrament. You cannot brand the breath of God.

The real path to greatness—the kind Kierkegaard called the "Knight of Faith"—looks more like obscurity than fame. It looks like holding fast to what is unseen. Like trusting in absurd grace when everything around you says, "monetize or die."

So, no, this isn’t a takedown of Big Mike. It’s a letter to the soul he leaves behind.

The Sacred Unbranding

Grace says you don’t have to be spectacular to be beloved. You don’t have to be consumable to be worthy. Your story is not your sales pitch. Your soul is not a statistic. You are not a brand. You are a human being, created in the image of a God who flipped the script on greatness by washing feet, not winning followers.

The world says, “Build a platform.” The gospel says, “Carry a cross.”

And that, dear reader, is the absurdity that might just save you.

Light a candle. Step out of the feed. Reclaim the holy ridiculousness of unmarketable grace.

Because real greatness can’t be sold. It must be lived.

\chapter{The Algorithm Is My Shepherd}

\section{Clickbait Theology and the End of Reverence}

"For they loved the praise of men more than the praise of God." — John 12:43 (KJV)

In the digital age, the age-old tension between truth and popularity has reached a fever pitch. The church, traditionally a place of sacred reverence, has increasingly become a platform where the pursuit of clicks and likes trumps the pursuit of godly wisdom. The Gospel, once a message of profound depth, has been reduced to a series of viral soundbites, designed for maximum engagement rather than transformation. This shift marks the emergence of what can only be described as clickbait theology—a form of ministry where the sacred is marketed for social media success.

In the same way that the prosperity gospel promised material wealth in exchange for faith, clickbait theology promises instant, tangible rewards—such as followers, shares, and online fame—in exchange for flashy, sensationalized messages that often lack the depth of Scripture and the gravity of Christ’s teachings. But where prosperity gospel preached that your faith would guarantee wealth, clickbait theology teaches that your faith can guarantee visibility. It's not about the eternal riches of the Kingdom of God; it's about the fleeting gratification of the “like” button.

\section{The Rise of the Digital Prosperity Gospel}

At its core, clickbait theology mirrors the same principles as the prosperity gospel, albeit in a more modern and socially acceptable form. Rather than teaching that God’s favor manifests in wealth, this new form of gospel teaches that God’s favor is demonstrated through platforms. If you’re not seen, you’re not blessed. If you’re not followed, you’re not succeeding.

Where once pastors preached for spiritual transformation, now they often preach for clicks. Sermons are no longer crafted to inspire deep, transformative reflection on the nature of God’s kingdom, but to provoke emotional reactions, designed for viral engagement. The message of salvation is reduced to a catchy headline, neatly packaged for the algorithm. What does this mean for the future of the Church?

The Bible speaks clearly about the dangers of seeking the praise of men over the praise of God, and John 12:43 says, “For they loved the praise of men more than the praise of God.” Yet in the digital age, this principle is reversed. The praise of men is now the goal; spiritual depth is a secondary concern, if a concern at all. The church has traded depth for attention, reverence for relevance.

\section{The Mechanization of Ministry}

The shift toward clickbait theology is no accident. It is the inevitable result of a world increasingly dominated by algorithms. Algorithms shape everything: from what we watch on YouTube to the content we consume on social media. They are the unseen shepherds guiding our thoughts, tastes, and even our spiritual beliefs. It is no longer enough for a sermon to be true—it must also be optimized for clicks. What was once a sacred responsibility has become a marketing campaign.

This is not just an issue of technology; it is an issue of theology. The church has been forced to adapt to the digital world, where success is measured by engagement metrics rather than spiritual fruit. Algorithms, by their very nature, do not care about the soul’s salvation. They care only about attention. And so, ministry has been reduced to a contest for visibility, where the message of the Gospel is bent to fit the expectations of a click-hungry world.

\section{The Paradox of Faith in a Click-Driven World}

Søren Kierkegaard, the great existentialist theologian, would have seen the dangers of this trend clearly. For Kierkegaard, faith was never about conforming to the expectations of the crowd. In his works, he constantly emphasized the necessity of standing alone before God in what he called the “leap of faith.” Faith, for Kierkegaard, was about existential commitment, not about pleasing others or catering to social approval.

In the context of today’s church, Kierkegaard’s insights take on a fresh relevance. When ministry is reduced to a popularity contest, it no longer serves its true purpose. If the faith is about catering to the masses, it becomes a faith of conformity, not conviction. The genuine leap of faith becomes impossible when the goal is not salvation, but social media validation. Kierkegaard’s “Knight of Faith,” the individual who chooses to follow God despite societal pressures, would be a rare figure in today’s church. In the world of clickbait theology, it’s the loudest voice, not the truest, that gets heard.

\section{The Algorithm’s Effect on the Soul}

While algorithms may seem neutral, they are far from it. Algorithms are designed to reward sensationalism and engagement, and in doing so, they reward content that is shallow, emotionally charged, and often devoid of substance. This is why the algorithm cannot be a shepherd. It does not lead us toward spiritual transformation or deeper communion with God—it leads us toward ever-increasing clicks.

The very nature of the algorithm encourages a shortcut spirituality, one where depth is sacrificed for instant gratification. Instead of encouraging deep study of Scripture or contemplation of divine mysteries, it encourages soundbites that provoke emotional reactions. It promotes viral content over sacred truth. And in this, it reduces faith to a transactional relationship—like for a like, share for a share—where the sacred becomes a commodity and the Gospel is merely another product to be consumed.

This stands in stark contrast to the teachings of Jesus, who repeatedly emphasized the need for deep, personal engagement with God. Matthew 6:6 speaks of the importance of private, personal prayer: “But thou, when thou prayest, enter into thy closet, and when thou hast shut thy door, pray to thy Father which is in secret; and thy Father which seeth in secret shall reward thee openly.” The algorithm rewards public displays of faith, while Jesus calls for private devotion. This contrast speaks volumes about the difference between genuine spiritual engagement and performative religiosity.

\section{The Recovery of Reverence}

So, what’s the way forward? How can we recover reverence in an age that demands spectacle? The answer is not necessarily to abandon digital platforms but to reclaim the sacred nature of ministry within them. The church must resist the pull of the algorithm and, instead, use technology as a tool for deeper engagement with God’s truth, not a mechanism for crowd-pleasing.

Charles Spurgeon, the Prince of Preachers, provides a model for what it looks like to preach with reverence in a world obsessed with numbers. Spurgeon’s ministry was marked by a commitment to truth, regardless of how it was received. He did not craft his sermons to please crowds or to gain followers; he crafted them to speak to the soul. In contrast to today’s performative preachers, Spurgeon’s sermons were filled with theological richness, designed not for virality but for spiritual transformation.

Spurgeon understood that ministry was not about securing applause or likes; it was about pointing people toward Christ. And this is where the church must return—away from the spectacle and toward the sacred. We must recover the slow, patient work of preaching the Gospel, regardless of how it performs on digital platforms.

\section{The End of the Algorithmic Shepherd}

In the end, the algorithm cannot shepherd us toward truth, and it cannot shepherd us toward God. It can only shepherd us toward what is popular, what is engaging, and what generates clicks. But the Christian calling is not about popularity. It is about faithfulness. It is about truth, and it is about love. The algorithm offers neither.

As the church navigates the digital age, it must resist the temptation to trade reverence for engagement. The soul cannot be reduced to a series of metrics or to a viral post. The soul longs for something deeper, something eternal. Only the true Shepherd can guide us there.


%%%%%%%%%%%%%%%%%%%%%%%%%%%%%%%%%%%%%%%%%%%%%%%%%%%%%%%%%%%%%%%%%%%%%%%%
% Auto-generated table of contents, list of figures and list of tables %
%%%%%%%%%%%%%%%%%%%%%%%%%%%%%%%%%%%%%%%%%%%%%%%%%%%%%%%%%%%%%%%%%%%%%%%%

\listoffigures
\listoftables


\end{document}